% Section 27: The cohomological meaning of the residue} | Part: 6 | Depends on: see docs/REWRITE_PLAN.md | Status: todo | Original lines: 3845-3854
%!TEX root = ../../../main.tex
\subsection{The cohomological meaning of the residue}\label{sec:rescohom}

For $X=\C\setminus\{a_1,\dots,a_k\}$ one has $H^1_{\mathrm{dR}}(X;\C)\cong\C^k$, generated by the classes of $\dd z/(z-a_j)$, and $H_1(X;\Z)\cong\Z^k$, generated by small circles around the $a_j$. The residue is the de Rham pairing between these:
$\Res_{a_j}f=\frac1{2\pi i}\bigl\langle[\text{circle around }a_j],[f\,\dd z]\bigr\rangle$.
A closed loop $\gamma$ in $X$ is homologous to $\sum_jn_j[\text{circle}_j]$, where $n_j=\frac1{2\pi i}\oint_\gamma\frac{\dd z}{z-a_j}$ is the winding number of $\gamma$ around $a_j$ (Example~\ref{ex:angle}); \eg a loop that goes twice around $a_j$ has $\oint\frac{\dd z}{z-a_j}=2\cdot2\pi i$, so $n_j=2$. This yields the general version
\begin{equation}\label{eq:winding}
\oint_\gamma f\,\dd z=2\pi i\sum_jn(\gamma,a_j)\Res_{a_j}f .
\end{equation}
The integral depends only on the homology class of the contour: a contour can be deformed freely as long as it does not cross a singularity, which is the practical principle behind contour integration.
